\chapter{Breast Cysts}
\index{Breast Cysts}

\section*{Definition and Overview}
A breast cyst is a fluid-filled sac that develops within the breast tissue, usually within the milk glands or ducts. Cysts are a common benign (non-cancerous) finding and represent part of the normal changes breast tissue undergoes in response to the hormones of the menstrual cycle. They may be single or multiple, affect one or both breasts, and range in size from a few millimetres to several centimetres. Simple breast cysts are not a form of cancer and do not become cancer, but they must be distinguished from other breast lumps -- including cancer -- by examination and imaging. A minority of cysts have more complex features on ultrasound, and these require additional investigation because of a small association with abnormal cells.

\section*{Epidemiology}
Breast cysts are very common, particularly in women in their 30s to 50s, around the perimenopause, when hormonal activity is at its most variable. They are rare after the menopause unless a woman is taking hormone replacement therapy. Many women have breast cysts without ever being aware of them, since cysts are frequently found incidentally on imaging performed for another reason or during screening. Cysts tend to fluctuate with the menstrual cycle, enlarging and becoming more tender before a period and shrinking or resolving afterwards, and they often recur over many years.

\section*{Aetiology and Pathophysiology}
Breast cysts arise from the glandular elements of the breast. As hormone levels fluctuate during the menstrual cycle, some ducts and lobules become distended and fill with fluid, forming a macroscopically visible cyst. This process is part of the broader pattern of \textbf{fibrocystic change}, in which breast tissue becomes lumpy, dense, and tender. Cysts can enlarge in the second half of the cycle, when hormonal stimulation is greatest, and shrink after menstruation. Simple cysts are smooth-walled sacs of fluid. A \textbf{complex cyst} contains internal echoes, septations, or solid-looking components on ultrasound, and although most prove benign, such cysts carry a small risk of malignancy and are investigated more thoroughly.

\section*{Clinical Features}
\begin{itemize}
    \item A smooth, round, or oval lump that feels firm or soft and often moves freely under the skin
    \item A lump that enlarges and becomes tender before a period and shrinks afterwards
    \item Breast pain or tenderness, sometimes of sudden onset
    \item Multiple lumps in one or both breasts, often accompanied by general lumpiness
    \item A lump discovered incidentally on imaging with no symptoms at all
    \item Occasional sudden enlargement with pain, which can follow bleeding into a cyst
\end{itemize}

\section*{Differential Diagnosis}
\begin{itemize}
    \item Fibroadenoma, a benign solid lump of glandular and connective tissue
    \item Breast cancer, which can occasionally feel indistinguishable from a cyst
    \item Fibrocystic change without a discrete cyst
    \item Galactocele, a milk-filled cyst in breastfeeding women
    \item Fat necrosis or a lipoma, benign fatty lesions
    \item Breast abscess, an infected and painful fluid collection
    \item Phyllodes tumour, a rare tumour that is usually benign
\end{itemize}

\section*{When to Seek Medical Care}
See a doctor about any new breast lump, especially one that persists through a full menstrual cycle or appears after the menopause. New lumps after the menopause warrant prompt investigation, because cysts are uncommon then and breast cancer becomes more likely. Any change in a previously known lump also requires review.

\textbf{Contact your local emergency services for:}
\begin{itemize}
    \item A breast lump accompanied by high fever, spreading redness, or severe pain -- possible infection or abscess
    \item A rapidly enlarging, hard, or fixed lump with skin changes such as dimpling or nipple retraction
    \item Bleeding from the nipple or a non-healing sore on the breast
\end{itemize}

\section*{Diagnosis and Investigations}
\begin{itemize}
    \item \textbf{Clinical examination:} the doctor assesses the lump's size, shape, consistency, and mobility, and examines the armpits and neck
    \item \textbf{Ultrasound:} the key investigation for a lump in a younger woman; a simple, smooth-walled, fluid-filled cyst is confidently diagnosed on ultrasound
    \item \textbf{Mammography:} used in women over about 35--40 years or when ultrasound is inconclusive, to exclude other lesions in surrounding tissue
    \item \textbf{Aspiration:} drainage of the cyst with a fine needle; if clear fluid is obtained and the lump disappears completely, no further treatment is needed
    \item \textbf{Biopsy:} if the aspirated fluid is bloodstained, the lump does not collapse fully, or the cyst looks complex on imaging, tissue sampling is performed to exclude cancer
    \item \textbf{Follow-up imaging:} repeat ultrasound or mammography to document resolution or monitor a complex cyst
\end{itemize}

\section*{Management}
\begin{itemize}
    \item \textbf{Reassurance and observation:} simple cysts that cause no symptoms can be left alone and usually need no treatment at all
    \item \textbf{Aspiration:} a painful or large cyst can be drained with a needle, which usually relieves discomfort immediately
    \item \textbf{Complex cysts:} these require aspiration and, where indicated, biopsy, with imaging follow-up to confirm stability
    \item \textbf{Pain relief:} simple analgesics and a well-fitting, supportive bra for cyclical tenderness
    \item \textbf{Follow-up imaging:} repeat ultrasound after some months if advised, to confirm the cyst has not recurred or changed
    \item \textbf{Continuing breast screening:} having cysts does not remove the need for routine mammography screening
\end{itemize}

\section*{Complications}
\begin{itemize}
    \item Recurrence of cysts after aspiration, which is common
    \item Infection within a cyst (rare), causing pain, redness, and fever
    \item Bleeding into a cyst, which can make it suddenly enlarge and become painful
    \item A small proportion of complex cysts prove to harbour abnormal cells or cancer, which is why they are investigated rather than dismissed
    \item Anxiety about a breast lump before it is confirmed to be benign
\end{itemize}

\section*{Prognosis and Outcomes}
Simple breast cysts are benign and do not increase the risk of breast cancer; they may come and go for years and usually resolve after the menopause. Complex cysts are investigated thoroughly because of a small associated risk, and women who have had them are typically followed up with imaging until stability is confirmed. With appropriate assessment, the vast majority of breast cysts prove harmless and require no more than reassurance, with occasional aspiration for comfort.

\section*{Prevention and Self-Care}
\begin{itemize}
    \item Know the normal look and feel of your breasts and report any new lump or change promptly
    \item Attend routine breast screening as recommended for your age group
    \item Wear a supportive, well-fitting bra, particularly when the breasts are tender
    \item Use simple pain relief and local warmth for discomfort if advised
    \item Return for scheduled follow-up imaging if a complex cyst was found
    \item Seek review if a known cyst enlarges, becomes painful, or otherwise changes
\end{itemize}

\section*{Sources and Further Reading}
The following organisations provide reliable, evidence-based information for patients, students, and clinicians.

\begin{itemize}
    \item NHS. \textit{Breast cysts: symptoms, diagnosis, and treatment information.} https://www.nhs.uk/conditions/
    \item Mayo Clinic. \textit{Breast cysts: overview and patient education.} https://www.mayoclinic.org/diseases-conditions/
    \item MSD Manual. \textit{Breast cysts and fibrocystic change: clinical reference.} https://www.msdmanuals.com/
    \item MedlinePlus. \textit{Breast cysts: consumer health information.} https://medlineplus.gov/
    \item National Institute for Health and Care Excellence. \textit{Clinical guidance on the assessment of breast lumps.} https://www.nice.org.uk/
    \item Centers for Disease Control and Prevention. \textit{Breast health and screening information.} https://www.cdc.gov/
\end{itemize}
